\documentclass[12pt]{article}
\newcommand{\bottomMargin}{2cm}

\newcommand{\header}{
	ДИЗАЙН И АНАЛИЗ НА АЛГОРИТМИ - \\ ПРАКТИКУМ\\
	\vspace{0.1cm}
	Летен семестър, 2026 г., трето контролно\\
	\vspace{0.1cm}
}
\newcommand{\tl}{$0.1$ сек.}
\newcommand{\ml}{$256$ MB}

\input{structure.tex}
\usepackage{newunicodechar}
\newunicodechar{ѝ}{{$\grave{\text{и}}$}}
\newunicodechar{Ѝ}{{$\grave{\text{И}}$}}

\begin{document}
	
\section{Задача К10. МАТРИЦИ}

Разглеждаме произведението на матрици $M_1 \times M_2 \times \ldots \times M_n$ с размерности $R_0 \times R_1$, $R_1 \times R_2$, $R_2 \times R_3, \ldots, R_{n-1} \times R_n$. Резултатът ще бъде един и същ от свойството за асоциативност, но броят на умноженията на реалните числа за получаването на този резултат може съществено да се различава. Да разгледаме следния пример:

Трябва да се намери произведението $M_1 \times M_2 \times M_3 \times M_4$, където матриците $M_1$, $M_2$, $M_3$ и $M_4$ са съответно с размерности $10 \times 20$, $20 \times 50$, $50 \times 1$ и $1 \times 100$. Ако пресмятаме произведението в реда, определен от скоби, по следния начин $M_1 \times (M_2 \times (M_3 \times M_4))$, то ще получим следния брой умножения на реални числа:

\begin{itemize}
	\item $50 \cdot 1 \cdot 100 = 5000$ -- от $(M_3 \times M_4)$ (получава се матрица с размерност $50 \times 100$)
	\item $20 \cdot 50 \cdot 100 = 100000$ -- от $M_2 \times (M_3 \times M_4)$ (получава се матрица с размерност $20 \times 100$)
	\item $10 \cdot 20 \cdot 100 = 20000$ -- от $M_1 \times (M_2 \times (M_3 \times M_4))$ (получава се матрица $10 \times 100$)
\end{itemize}

При тази последователност на умножаване на матриците са нужни $125\,000$ умножения на реални числа. Нека да извършим умножението в друг ред: $(M_1 \times (M_2 \times M_3)) \times M_4$. Броят умножения на реални числа за получаване на същия резултат ще бъде:

\begin{itemize}
	\item $20 \cdot 50 \cdot 1 = 1000$ -- от $(M_2 \times M_3)$ (получава се матрица с размерност $20 \times 1$)
	\item $10 \cdot 20 \cdot 1 = 200$ -- от $(M_1 \times (M_2 \times M_3))$ (получава се матрица с размерност $10 \times 1$)
	\item $10 \cdot 1 \cdot 100 = 1000$ -- от $(M_1 \times (M_2 \times M_3)) \times M_4$ (получава се матрица с размерност $10 \times 100$)
\end{itemize}

При тази последователност на умножаване на матриците са нужни $2200$ умножения на реални числа -- разликата е чувствителна, нали?

Напишете програма \textbf{\texttt{matrices}}, която намира минималния брой умножения на реални числа, които са необходими, за да бъде намерено произведението на $n$ матрици $M_1 \times M_2 \times \ldots \times M_n$ с размерности $R_0 \times R_1$, $R_1 \times R_2$, $R_2 \times R_3, \ldots, R_{n-1} \times R_n$.

\subsection{Вход}
От първия ред на стандартния вход се въвежда цяло положително число $n$ -- брой на матриците. От втория ред се въвеждат $n+1$ цели положителни числа $R_0, R_1, R_2, \ldots, R_n$, разделени с интервали, които задават размерностите на матриците.

\subsection{Изход}
На един ред на стандартния изход изведете само едно число -- намерения минимален брой умножения на реални числа, необходими за намирането на произведението на матриците.

\subsection{Ограничения}
\vspace{0.1em}
\begin{itemize}[label=$\clubsuit$]
	\item $2 \leq n \leq 100$
	\item $1 \leq R_0, R_1, R_2, \ldots, R_n \leq 1000$
	\item В 30\% от тестовете $2 \leq n \leq 10$
\end{itemize}

\subsection{Пример}
\begin{table}[!ht]
	\begin{tblr}{|X[50,l]|X[50,l]|}
		\hline
		\textbf{Вход} & \textbf{Изход} \\
		\hline
		\texttt{4 \\
			10 20 50 1 100}
		&
		\texttt{2200}
		\\
		\hline
	\end{tblr}
	\caption*{Примерът съответства на примера, разгледан в условието.}
\end{table}
\FloatBarrier

\end{document}
